% The bench for this chapter, and the two transports in build order.
\begin{tikzpicture}[font=\sffamily\scriptsize, x=1cm, y=1cm]

% ================================================================ the boards
\node[board, minimum width=40mm, minimum height=24mm] (nuc) at (-4.0,0.6) {NUCLEO\\H7A3ZI-Q};
\node[board, fill=red!35!black!60, minimum width=36mm, minimum height=24mm] (pi) at (4.0,0.6)
  {Raspberry Pi 4\\the agent};

% ================================================================ transport 1: the serial port
\draw[usbcable] (nuc.north) -- (-4.0,3.0) -- (4.0,3.0) -- (pi.north);
\node[buslabel, anchor=south] at (0,3.05) {1. the probe's serial port, 921600 baud, transfers both ways};
\node[chlabel, anchor=north] at (0,2.95) {shipped, proven, already wired: the milestone that proves the port};

% ================================================================ transport 2: the bus
\node[hw, text width=18mm, minimum height=8mm] (tx1) at (-1.6,-1.2) {transceiver};
\node[hw, text width=18mm, minimum height=8mm] (tx2) at (1.6,-1.2) {adapter};
\draw[cable] (nuc.south) -- (-4.0,-1.2) -- (tx1.west);
\draw[busline] (tx1.east) -- (tx2.west);
\draw[cable] (tx2.east) -- (4.0,-1.2) -- (pi.south);
\node[buslabel, anchor=north] at (0,-1.75) {2. the joint bus, flexible-data, 64-byte payload};
\node[chlabel, anchor=north] at (0,-2.15) {not a shipped transport: six functions, written here};

% ================================================================ the build machine
\node[ext, text width=30mm, minimum height=11mm] (lap) at (-4.6,-3.3)
  {the authoring laptop:\\the container build};
\draw[flow, dashed] (lap.west) -- (-6.9,-3.3) -- (-6.9,0.6) -- (nuc.west);
\node[chlabel, anchor=west, text width=22mm] at (-6.8,-1.4)
  {a static library, built there and copied here};

% ================================================================ what is not here
\node[absentblk, text width=28mm, minimum height=11mm] (eth) at (-0.6,-3.3)
  {no Ethernet controller on this part};
\node[laterblk, text width=28mm, minimum height=11mm] (usb) at (3.4,-3.3)
  {a user USB connector, unconfirmed};

% ================================================================ captions
\node[chlabel, anchor=north, text width=32mm, align=center] at (-4.6,-4.0)
  {the image does not state which architectures it supports, so the build runs
   where the answer is certain};
\node[chlabel, anchor=north, text width=32mm, align=center] at (-0.6,-4.0)
  {which removes the UDP transport before it is considered, on evidence from
   chapter 4 rather than from a product page};
\node[chlabel, anchor=north, text width=32mm, align=center] at (3.4,-4.0)
  {the USB serial transport depends on reading the silkscreen, so this chapter
   does not assume it};

\node[umlnote, text width=112mm, anchor=north west] at (-6.4,-5.4)
  {Two transports, built in that order, for the same reason chapter 9 started in
   loopback: the first one proves the port with a known-good path, so that when the
   second one fails the fault can only be in the transport. Reversing the order
   turns two unknowns into one long evening.};
\end{tikzpicture}
